\chapter{Persamaan Garis dan Sistem Persamaan Linear}\label{ch:g10:lines}

Sebuah garis lurus pada bidang \omterm{def:g10:coordgeom:system}{koordinat} diperikan oleh sebuah \omterm{def:g10:algebra:equation}{persamaan}, dan
dua garis bersama-sama membentuk sistem \omterm{def:g10:algebra:equation}{persamaan} yang penyelesaiannya adalah
titik potongnya. Bab ini menghubungkan tiga bahasa --- geometris
(garis), aljabar (\omterm{def:g10:algebra:equation}{persamaan}) dan fungsional (\omterm{def:g10:reffunc:affine}{fungsi afin} dari
\cref{ch:g10:reffunc}) --- serta mengajarkan dua cara baku menyelesaikan
\omterm{def:g10:lines:system}{sistem persamaan linear}.

\section{Persamaan sebuah garis}

\begin{theorem}[Persamaan tereduksi]\label{thm:g10:lines:reduced}
Pada sebuah \omterm{def:g10:coordgeom:system}{sistem koordinat}:
\begin{enumerate}
  \item setiap garis yang tidak tegak mempunyai \omterm{def:g10:algebra:equation}{persamaan} berbentuk
        \[
        y = mx + p ,
        \]
        dengan $m$ (\emph{gradien}\index{gradien} garisnya) dan $p$
        (\emph{titik potong sumbu-$y$}) tertentu secara tunggal --- inilah
        \emph{persamaan tereduksi}\index{persamaan tereduksi} garis itu;
  \item setiap garis tegak mempunyai \omterm{def:g10:algebra:equation}{persamaan} $x = c$;
  \item sebuah titik terletak pada garis itu tepat ketika \omterm{def:g10:coordgeom:system}{koordinatnya}
        memenuhi \omterm{def:g10:algebra:equation}{persamaannya}.
\end{enumerate}
\end{theorem}
\admitted

\begin{proposition}[Gradien lewat dua titik]\label{prop:g10:lines:slope}
\omterm{def:g10:reffunc:affine}{Gradien} garis (yang tidak tegak) melalui $A(x_A, y_A)$ dan
$B(x_B, y_B)$, dengan $x_A \neq x_B$, adalah
\[
m = \frac{y_B - y_A}{x_B - x_A}
\qquad
\text{(``perubahan $y$ dibagi perubahan $x$'').}
\]
\end{proposition}

\begin{proof}
Kedua titik memenuhi $y = mx + p$, jadi $y_B - y_A = (mx_B + p) - (mx_A + p)
= m(x_B - x_A)$; bagilah dengan $x_B - x_A \neq 0$.
\end{proof}

\begin{omfigure}
\begin{tikzpicture}
\begin{axis}[omaxis,
  width=8.2cm, height=5.8cm,
  xmin=-1.4, xmax=4.4, ymin=-1.4, ymax=4.4,
  xtick={1,2,3}, ytick={1,3}]
  \addplot[omDef, thick, domain=-1.1:4.1] {x + 0.5};
  \addplot[omThm, thick, domain=-1.1:4.1] {3};
  \draw[omMeth, thick] (axis cs:2,-1.2) -- (axis cs:2,4.2);
  \node[omDef, font=\small, rotate=38] at (axis cs:2.7,3.6)
    {$y = x + 0.5$};
  \node[omThm, font=\small] at (axis cs:-0.55,3.3) {$y = 3$};
  \node[omMeth, font=\small, right] at (axis cs:2,-0.8) {$x = 2$};
\end{axis}
\end{tikzpicture}

{\small Tiga jenis garis: miring ($m \neq 0$), mendatar ($m = 0$),
dan tegak (tanpa \omterm{thm:g10:lines:reduced}{persamaan tereduksi} --- \omterm{def:g10:algebra:equation}{persamaannya} $x = c$).}
\end{omfigure}

\begin{method}[Mencari persamaan garis lewat dua titik]\label{met:g10:lines:through}
Diberikan $A$ dan $B$ dengan $x_A \neq x_B$:
\begin{enumerate}
  \item hitunglah \omterm{def:g10:reffunc:affine}{gradiennya} $m = \dfrac{y_B - y_A}{x_B - x_A}$;
  \item tulislah $y = mx + p$ lalu substitusikan \omterm{def:g10:coordgeom:system}{koordinat} $A$ (atau $B$)
        untuk mencari $p$;
  \item periksalah \omterm{def:g10:algebra:equation}{persamaan} akhirnya dengan titik yang \emph{satunya}.
\end{enumerate}
Jika $x_A = x_B$, garisnya tegak: \omterm{def:g10:algebra:equation}{persamaannya} $x = x_A$.
\end{method}

\begin{example}\label{ex:g10:lines:through}
Garis melalui $A(1, 5)$ dan $B(3, 1)$. \omterm{def:g10:reffunc:affine}{Gradiennya}:
$m = \frac{1 - 5}{3 - 1} = \frac{-4}{2} = -2$. Lalu $y = -2x + p$; dari
titik $A$ diperoleh $5 = -2 \times 1 + p$, jadi $p = 7$. \omterm{def:g10:algebra:equation}{Persamaannya}:
$y = -2x + 7$. Periksa dengan $B$: $-2 \times 3 + 7 = 1$.
\end{example}

\begin{proposition}[Garis sejajar]\label{prop:g10:lines:parallel}
Dua garis yang tidak tegak sejajar jika dan hanya jika \omterm{def:g10:reffunc:affine}{gradiennya} sama.
Dua garis sejajar yang berbeda tidak punya titik persekutuan; dua garis
dengan \omterm{def:g10:reffunc:affine}{gradien} berbeda punya tepat satu.
\end{proposition}

\begin{proof}
Misalkan $y = mx + p$ dan $y = m'x + p'$ kedua garis itu. Titik persekutuan
memenuhi $mx + p = m'x + p'$, yaitu $(m - m')x = p' - p$. Jika $m \neq m'$,
\omterm{def:g10:algebra:equation}{persamaan} ini punya tepat satu penyelesaian $x = \frac{p' - p}{m - m'}$: satu
titik potong. Jika $m = m'$, \omterm{def:g10:algebra:equation}{persamaannya} berbunyi $0 = p' - p$: tidak ada
penyelesaian bila $p \neq p'$ (dua sejajar yang berbeda), dan setiap $x$ bila
$p = p'$ (garis yang sama). Jadi kesejajaran (tidak berpotongan, atau berimpit)
setara dengan $m = m'$.
\end{proof}

\begin{remark}[Vektor arah]
Garis $y = mx + p$ naik $m$ satuan untuk setiap satu satuan ke kanan, sehingga
\omterm{def:g10:vectors:vector}{vektor} $\vec u\,(1, m)$ adalah sebuah \emph{vektor
arah}\index{vektor arah} garis itu: garisnya sejajar dengan
$\vec u$. Dua garis sejajar tepat ketika \omterm{def:g10:vectors:vector}{vektor} arahnya
\omterm{def:g10:vectors:collinear}{kolinear} (\cref{thm:g10:vectors:det}).
\end{remark}

\section{Sistem persamaan linear}

\begin{definition}[Sistem persamaan linear]\label{def:g10:lines:system}
Sebuah \emph{sistem dua persamaan linear}\index{sistem persamaan linear} dalam
variabel $x$ dan $y$ adalah sepasang \omterm{def:g10:algebra:equation}{persamaan}
\[
\begin{cases}
a x + b y = c \\
a' x + b' y = c'
\end{cases}
\]
yang harus dipenuhi \emph{serentak}. Sebuah \emph{penyelesaian} adalah pasangan
$(x, y)$ yang memenuhi keduanya. Setiap \omterm{def:g10:algebra:equation}{persamaannya} memerikan sebuah garis,
sehingga menyelesaikan sistem itu berarti mencari perpotongan dua garis.
\end{definition}

\begin{method}[Substitusi]\label{met:g10:lines:substitution}
\begin{enumerate}
  \item Selesaikan salah satu \omterm{def:g10:algebra:equation}{persamaan} untuk salah satu variabel --- pilih
        yang paling mudah, misalnya nyatakan $y$ dalam $x$;
  \item substitusikan bentuk ini ke \omterm{def:g10:algebra:equation}{persamaan} yang \emph{satunya}, yang kini
        hanya memuat $x$;
  \item selesaikan untuk $x$, lalu hitunglah $y$ dari langkah 1;
  \item periksalah pasangan itu pada kedua \omterm{def:g10:algebra:equation}{persamaan} semula.
\end{enumerate}
\end{method}

\begin{example}\label{ex:g10:lines:substitution}
Selesaikan $\begin{cases} x + 2y = 8 \\ 3x - y = 3 \end{cases}$ dengan
substitusi. \omterm{def:g10:algebra:equation}{Persamaan} kedua memberi $y = 3x - 3$. Substitusikan ke yang
pertama:
\[
x + 2(3x - 3) = 8
\ \Longleftrightarrow\
x + 6x - 6 = 8
\ \Longleftrightarrow\
7x = 14
\ \Longleftrightarrow\
x = 2 ,
\]
lalu $y = 3 \times 2 - 3 = 3$. Penyelesaiannya adalah pasangan $(2, 3)$.
Periksa: $2 + 2 \times 3 = 8$ dan $3 \times 2 - 3 = 3$.
\end{example}

\begin{method}[Eliminasi]\label{met:g10:lines:elimination}
\begin{enumerate}
  \item Kalikan setiap \omterm{def:g10:algebra:equation}{persamaan} dengan bilangan yang dipilih baik-baik agar
        salah satu variabelnya berkoefisien saling berlawanan pada kedua \omterm{def:g10:algebra:equation}{persamaan};
  \item jumlahkan kedua \omterm{def:g10:algebra:equation}{persamaan}: variabel itu lenyap;
  \item selesaikan \omterm{def:g10:algebra:equation}{persamaan} satu variabel yang tersisa, lalu substitusikan
        kembali untuk mencari variabel yang lain;
  \item periksalah pasangan itu pada kedua \omterm{def:g10:algebra:equation}{persamaan} semula.
\end{enumerate}
\end{method}

\begin{example}\label{ex:g10:lines:elimination}
Selesaikan $\begin{cases} 2x + 3y = 7 \\ 5x - 2y = 8 \end{cases}$ dengan
eliminasi. Untuk melenyapkan $y$, kalikan \omterm{def:g10:algebra:equation}{persamaan} pertama dengan $2$ dan
yang kedua dengan $3$:
\[
\begin{cases}
4x + 6y = 14 \\
15x - 6y = 24
\end{cases}
\]
Menjumlahkan keduanya: $19x = 38$, jadi $x = 2$. Substitusi ke $2x + 3y = 7$:
$4 + 3y = 7$, jadi $y = 1$. Penyelesaiannya: $(2, 1)$. Periksa pada \omterm{def:g10:algebra:equation}{persamaan}
kedua: $5 \times 2 - 2 \times 1 = 8$.
\end{example}

\begin{omfigure}
\begin{tikzpicture}
\begin{axis}[omaxis,
  width=8.2cm, height=6.0cm,
  xmin=-0.9, xmax=4.9, ymin=-1.4, ymax=4.9,
  xtick={2}, ytick={1}]
  \addplot[omDef, thick, domain=-0.6:4.6] {(7 - 2*x)/3};
  \addplot[omThm, thick, domain=-0.5:2.9] {(5*x - 8)/2};
  \fill (axis cs:2,1) circle (2pt);
  \node[above right, font=\small] at (axis cs:2,1) {$(2, 1)$};
  \node[omDef, font=\small] at (axis cs:4.05,0.35) {$2x + 3y = 7$};
  \node[omThm, font=\small] at (axis cs:1.3,3.3) {$5x - 2y = 8$};
\end{axis}
\end{tikzpicture}

{\small \omterm{def:g10:lines:system}{Sistem persamaan linear} yang dilihat secara geometris: setiap
\omterm{def:g10:algebra:equation}{persamaan} adalah sebuah garis, dan penyelesaian $(2,1)$ pada
\cref{ex:g10:lines:elimination} adalah titik potongnya.}
\end{omfigure}

\begin{remark}[Berapa banyak penyelesaiannya?]
Seperti dua garis, \omterm{def:g10:lines:system}{sistem persamaan linear} pada umumnya punya tepat satu
penyelesaian (garis tidak sejajar), dan pada keadaan khusus tidak punya
penyelesaian (dua sejajar yang berbeda) atau punya tak hingga banyak (dua kali
garis yang sama). Kriteria \cref{thm:g10:vectors:det} memutuskannya: sistem itu
punya penyelesaian tunggal tepat ketika $ab' - a'b \neq 0$.
\end{remark}

\begin{example}[Memodelkan dengan sistem]\label{ex:g10:lines:modeling}
Tiga buku tulis dan dua pena berharga $8$; satu buku tulis dan empat pena
berharga $6$. Misalkan $x$ harga sebuah buku tulis dan $y$ harga sebuah pena:
\[
\begin{cases}
3x + 2y = 8 \\
x + 4y = 6 .
\end{cases}
\]
Dari \omterm{def:g10:algebra:equation}{persamaan} kedua, $x = 6 - 4y$; setelah disubstitusikan,
$3(6 - 4y) + 2y = 8$, jadi $18 - 10y = 8$, yang memberi $y = 1$ lalu
$x = 2$. Sebuah buku tulis berharga $2$ dan sebuah pena berharga $1$.
\end{example}

\section{Latihan}

\begin{exercise}[$\star$]\label{exo:g10:lines:1}
Untuk setiap garis, bacalah \omterm{def:g10:reffunc:affine}{gradien} dan titik potong sumbu-$y$-nya, lalu sketsalah:
\[
y = 2x - 3, \qquad
y = -\tfrac13 x + 1, \qquad
y = 4, \qquad
x = -2 .
\]
\end{exercise}

\begin{exercise}[$\star$]\label{exo:g10:lines:2}
Apakah titik $A(2, 7)$ terletak pada garis $y = 3x + 1$? Lalu $B(-1, -1)$?
Lalu titik $C(4, 13)$?
\end{exercise}

\begin{exercise}[$\star$]\label{exo:g10:lines:3}
Carilah \omterm{thm:g10:lines:reduced}{persamaan tereduksi} garis yang melalui:
\[
\text{(a) } A(0, 2) \text{ dan } B(3, 8);
\qquad
\text{(b) } A(-1, 5) \text{ dan } B(2, -4);
\qquad
\text{(c) } A(4, 1) \text{ dan } B(4, 6).
\]
\end{exercise}

\begin{exercise}[$\star$]\label{exo:g10:lines:4}
Di antara garis $y = 3x - 1$, $y = -3x + 2$, $y = 3x + 5$,
$y = \frac{6x + 2}{2}$, manakah yang saling sejajar? Manakah yang
sebenarnya berimpit?
\end{exercise}

\begin{exercise}[$\star$]\label{exo:g10:lines:5}
Selesaikan dengan substitusi:
\[
\begin{cases}
y = 2x - 1 \\
3x + y = 9
\end{cases}
\qquad\text{lalu}\qquad
\begin{cases}
x - y = 4 \\
2x + 3y = 3 .
\end{cases}
\]
\end{exercise}

\begin{exercise}[$\star$]\label{exo:g10:lines:6}
Selesaikan dengan eliminasi:
\[
\begin{cases}
x + y = 10 \\
x - y = 4
\end{cases}
\qquad\text{lalu}\qquad
\begin{cases}
3x + 2y = 12 \\
2x + 5y = 8 .
\end{cases}
\]
\end{exercise}

\begin{exercise}[$\star\star$]\label{exo:g10:lines:7}
Tentukan, tanpa menyelesaikannya, berapa penyelesaian yang dipunyai tiap sistem:
\[
\begin{cases}
2x - y = 3 \\
-4x + 2y = -6
\end{cases}
\qquad
\begin{cases}
2x - y = 3 \\
-4x + 2y = 1
\end{cases}
\qquad
\begin{cases}
2x - y = 3 \\
x + y = 0 .
\end{cases}
\]
\end{exercise}

\begin{exercise}[$\star\star$]\label{exo:g10:lines:8}
Dua ratus lembar karcis terjual untuk sebuah pementasan sekolah, sebagian
seharga $5$ (anak-anak) dan sebagian seharga $9$ (dewasa), dengan jumlah
$1352$. Berapa lembar karcis dari masing-masing jenis yang terjual?
\end{exercise}

\begin{exercise}[$\star\star$]\label{exo:g10:lines:9}
Misalkan $d$ garis $y = 2x - 3$ dan $A(1, 4)$.
\begin{enumerate}
  \item Berikan \omterm{def:g10:algebra:equation}{persamaan} garis $d'$ yang sejajar $d$ dan melalui
        $A$.
  \item Hitunglah titik potong $d'$ dengan sumbu $x$.
\end{enumerate}
\end{exercise}

\begin{exercise}[$\star\star\star$]\label{exo:g10:lines:10}
Misalkan $A(0, 0)$, $B(6, 2)$ dan $C(2, 4)$.
\begin{enumerate}
  \item Carilah \omterm{def:g10:algebra:equation}{persamaan} garis berat segitiga $ABC$ yang ditarik
        dari $A$ dan dari $B$ (sebuah garis berat menghubungkan titik sudut
        dengan \omterm{prop:g10:coordgeom:midpoint}{titik tengah} sisi di hadapannya).
  \item Hitunglah titik potongnya $G$, lalu periksa bahwa
        $G$ juga terletak pada garis berat ketiga. (Semestinya diperoleh bahwa
        \omterm{def:g10:coordgeom:system}{koordinat} $G$ adalah rata-rata \omterm{def:g10:coordgeom:system}{koordinat} $A$, $B$, $C$.)
\end{enumerate}
\end{exercise}

\section{Soal: Satu penyelesaian, tidak ada, atau tak hingga banyak}

\begin{problem}[{Soal akhir pekan --- dua garis mempunyai tiga nasib
yang mungkin, sebuah determinan memutuskan di antaranya, dan buku
pelajaran tertua di dunia sudah mengetahuinya}]
\label{pb:g10:lines:1}
Dua \omterm{def:g10:algebra:equation}{persamaan} linear, dua variabel: sepasang garis yang digambarnya
dapat berpotongan sekali, tidak pernah, atau justru berimpit --- dan
\emph{setiap} \omterm{def:g10:lines:system}{sistem persamaan linear} mewarisi salah satu dari ketiga nasib
itu. Soal ini menggolongkannya, menjumpai bilangan yang memutuskannya (seorang
kenalan lama dari bab \omterm{def:g10:vectors:vector}{vektor}), menyelesaikan soal burung pegar dan kelinci
berumur dua ribu tahun dari kitab Tiongkok
\emph{Sembilan Bab}, lalu berakhir pada titik sudut sebuah daerah --- yaitu
langkah pertama pengoptimuman yang dipakai setiap maskapai dan pabrik
dewasa ini.

\medskip
\noindent\textbf{Bagian I --- Garis, dengan fasih.}
\begin{enumerate}
  \item Carilah \omterm{thm:g10:lines:reduced}{persamaan tereduksi} garis yang melalui
        $(1, 2)$ dan $(3, 8)$
        (\cref{met:g10:lines:through}).
  \item Di antara $y = 3x - 1$, $y = 3x + 4$ dan $y = -x + 7$: dua
        garis yang mana yang sejajar
        (\cref{prop:g10:lines:parallel}), dan di mana pasangan yang tidak
        sejajar bertemu? (Hitunglah satu titik potong.)
  \item Garis melalui $(2, 1)$ dan $(2, 5)$: mengapa garis itu tidak
        mempunyai \omterm{thm:g10:lines:reduced}{persamaan tereduksi} $y = mx + p$, dan apa
        \omterm{def:g10:algebra:equation}{persamaannya}?
  \item Apakah titik $(4, 11)$ terletak pada garis $y = 3x - 1$? Lalu
        $(-1, -5)$?
  \item Garis melalui $(0, 3)$ dengan \omterm{def:g10:reffunc:affine}{gradien} $-2$: berikan
        \omterm{def:g10:algebra:equation}{persamaannya}, kedua titik potong sumbunya, dan luas
        segitiga yang dipotongnya dari kuadran pertama.
\end{enumerate}

\medskip
\noindent\textbf{Bagian II --- Tiga nasib.}
\begin{enumerate}[resume]
  \item Selesaikan dengan substitusi
        (\cref{met:g10:lines:substitution}):
        $\begin{cases} y = 2x - 3 \\ 3x + y = 7 \end{cases}$
  \item Selesaikan dengan eliminasi
        (\cref{met:g10:lines:elimination}):
        $\begin{cases} 2x + 3y = 8 \\ 5x - 3y = -1 \end{cases}$
  \item Golongkan --- lalu tafsirkan dengan garis --- kedua
        sistem berikut
        \[
        \begin{cases} 2x - y = 3 \\ 4x - 2y = 6 \end{cases}
        \qquad\text{dan}\qquad
        \begin{cases} 2x - y = 3 \\ 4x - 2y = 1 . \end{cases}
        \]
        Nyatakan trikotomi lengkapnya: berpotongan, sejajar,
        berimpit --- satu penyelesaian, tidak ada, tak hingga banyak.
  \item Untuk sistem umum $ax + by = e$, $cx + dy = f$,
        bilangan yang memutuskannya adalah $ad - bc$ --- determinan yang
        sama dengan uji kekolinearan pada
        \cref{pb:g10:vectors:1}. Jelaskan kebetulan itu (dua
        \omterm{def:g10:vectors:vector}{vektor} yang mana yang \omterm{def:g10:vectors:collinear}{kolinear} tepat ketika garisnya
        sejajar atau berimpit?), lalu hitunglah $ad - bc$ untuk
        ketiga sistem pada pertanyaan 7 dan 8.
  \item Untuk nilai $m$ yang mana
        $\begin{cases} x + my = 3 \\ 2x + 6y = 7 \end{cases}$
        tidak mempunyai penyelesaian? Selesaikan sistem itu untuk semua nilai
        $m$ yang lain\,\dots\ atau setidaknya jelaskan bagaimana caranya.
\end{enumerate}

\medskip
\noindent\textbf{Bagian III --- Kitab Sembilan Bab.}
\begin{enumerate}[resume]
  \item Dari \emph{Sembilan Bab tentang Seni Matematika}
        (Tiongkok, sekitar 2000 tahun silam): sebuah kandang berisi burung
        pegar dan kelinci --- seluruhnya $35$ kepala dan $94$ kaki. Berapa
        ekor masing-masing?
  \item Sebuah kedai jus mencampur sirop berkadar buah $60\,\%$ dengan
        minuman berkadar buah $20\,\%$ untuk memperoleh $30$ L berkadar
        buah $35\,\%$. Berapa liter masing-masing?
  \item Sebuah bilangan dua angka mempunyai jumlah angka $12$; menukar kedua
        angkanya mengurangi bilangan itu sebesar $18$. Carilah bilangan itu,
        dengan memakai aljabar angka pada soal akhir pekan tentang aljabar
        angka di jilid sebelumnya untuk menyusun sistemnya.
  \item Di toko roti, $3$ kopi dan $2$ croissant berharga
        $9.10$ euro; $2$ kopi dan $3$ croissant berharga
        $8.90$ euro. Selesaikan dengan anggun: apa yang langsung diberitahukan
        oleh \emph{menjumlahkan} kedua \omterm{def:g10:algebra:equation}{persamaan} dan oleh
        \emph{mengurangkannya}?
  \item Sebuah lapak pasar mengaku: $2$ apel $+ 1$ pir seharga
        $5$ euro, dan $4$ apel $+ 2$ pir seharga $9$ euro.
        Selesaikan --- atau lebih tepat, jelaskan apa yang disingkapkan
        nasib sistem itu tentang aritmetika si pedagang.
\end{enumerate}

\medskip
\noindent\textbf{Bagian IV --- Titik sudut yang menentukan.}
\begin{enumerate}[resume]
  \item Garis $y = 3x - 1$ membelah bidang menjadi dua. Di sisi
        mana titik asalnya berada? Perikan himpunan
        $y > 3x - 1$ dan bagaimana cara menguji sebuah titik terhadapnya.
  \item Gambarlah daerah yang ditentukan oleh keempat kendala
        \[
        x \geq 0, \qquad y \geq 0, \qquad x + y \leq 4,
        \qquad y \leq 2x + 1,
        \]
        lalu hitunglah keempat titik sudutnya.
  \item Laba sebuah bengkel adalah $P = 3x + 2y$ pada daerah di
        pertanyaan 17. Hitunglah $P$ di keempat titik sudutnya. Dengan
        menerima bahwa \omterm{def:g10:functions:extrema}{maksimum} besaran linear semacam itu pada daerah
        semacam itu selalu dicapai di sebuah titik sudut (bayangkan
        keluarga garis $3x + 2y$ tetap yang menyapu daerahnya), berikan
        rencana produksi yang optimum.
  \item Kedua taksi pada soal akhir pekan tentang dua termometer di jilid
        sebelumnya sebagai sebuah sistem: tulislah lalu selesaikan
        $\begin{cases} y = 1.5x + 3 \\ y = 2x \end{cases}$ lalu
        tafsirkan penyelesaiannya.
  \item Penutup, kamusnya: satu \omterm{def:g10:algebra:equation}{persamaan} $\leftrightarrow$
        satu garis; satu sistem $\leftrightarrow$ dua garis dengan
        tiga nasib, yang diputuskan oleh $ad - bc$; banyak pertidaksamaan
        $\leftrightarrow$ sebuah daerah yang titik sudutnya mengusung
        nilai optimum. Nyatakan itu dalam tiga kalimat --- kamu baru saja
        berkeliling serambi depan \emph{aljabar linear} dan
        \emph{pemrograman linear}, yang keduanya dibangun utuh pada
        jilid universitas.

\end{enumerate}
\end{problem}
